\textbf{Решение задачи 2-1}

\vspace{1em}
\textbf{1. Анализ исходных данных и расчет дисперсии шума}

Из условия задачи известны следующие параметры преселектора (БВЧ):
\begin{itemize}
    \item Шумовая полоса пропускания: $\Pi_{\text{ш}} = 1 \text{ МГц} = 10^6 \text{ Гц}$
    \item Коэффициент передачи проходной мощности: $K_{p\text{(дБ)}} = 6 \text{ дБ}$
    \item Спектральная плотность приведенного шума: $G_{\text{вх}} = 5 \cdot 10^{-19} \text{ Вт/Гц}$
    \item Сопротивление нагрузки: $R_{\text{н}} = 50 \text{ Ом}$
\end{itemize}

Переведем коэффициент передачи из децибел в разы. В радиотехнических расчетах $6 \text{ дБ}$ для мощности соответствует увеличению практически ровно в 4 раза:
\begin{equation*}
    K_p = 10^{\frac{6}{10}} \approx 3.98 \approx 4 \text{ раза}
\end{equation*}

Мощность шума на выходе преселектора рассчитывается как произведение спектральной плотности на полосу пропускания и коэффициент усиления:
\begin{equation*}
    P_{\text{вых}} = G_{\text{вх}} \cdot \Pi_{\text{ш}} \cdot K_p = 5 \cdot 10^{-19} \cdot 10^6 \cdot 4 = 20 \cdot 10^{-13} = 2 \cdot 10^{-12} \text{ Вт}
\end{equation*}

Определим дисперсию напряжения шума на выходе $\sigma^2$, зная сопротивление нагрузки:
\begin{equation*}
    \sigma^2 = P_{\text{вых}} \cdot R_{\text{н}} = 2 \cdot 10^{-12} \cdot 50 = 100 \cdot 10^{-12} = 10^{-10} \text{ В}^2
\end{equation*}

Эффективное напряжение (среднеквадратическое отклонение) шума:
\begin{equation*}
    \sigma = \sqrt{10^{-10}} = 10^{-5} \text{ В} = 10 \text{ мкВ}
\end{equation*}

\vspace{1em}
\textbf{2. Плотность вероятности мгновенных значений напряжения шума}

Мгновенные значения узкополосного шума на выходе БВЧ подчиняются нормальному (гауссовскому) закону распределения с нулевым математическим ожиданием. Аналитическое выражение:
\begin{equation*}
    W(u) = \frac{1}{\sigma \sqrt{2\pi}} \exp\left(-\frac{u^2}{2\sigma^2}\right)
\end{equation*}
Подставляя найденное значение $\sigma = 10^{-5} \text{ В}$:
\begin{equation*}
    W(u) = \frac{1}{10^{-5} \sqrt{2\pi}} \exp\left(-\frac{u^2}{2 \cdot 10^{-10}}\right) = \frac{10^5}{\sqrt{2\pi}} \exp\left(-5 \cdot 10^9 \cdot u^2\right) \quad [\text{В}^{-1}]
\end{equation*}

\vspace{1em}
\textbf{3. Плотность вероятности огибающей шума}

Огибающая узкополосного гауссовского шума на выходе БВЧ подчиняется закону распределения Релея. Аналитическое выражение:
\begin{equation*}
    W(E) = \frac{E}{\sigma^2} \exp\left(-\frac{E^2}{2\sigma^2}\right), \quad \text{при } E \ge 0
\end{equation*}
Подставляя дисперсию $\sigma^2 = 10^{-10} \text{ В}^2$:
\begin{equation*}
    W(E) = \frac{E}{10^{-10}} \exp\left(-\frac{E^2}{2 \cdot 10^{-10}}\right) = 10^{10} \cdot E \cdot \exp\left(-5 \cdot 10^9 \cdot E^2\right) \quad [\text{В}^{-1}]
\end{equation*}

\vspace{1em}
\textbf{4. Графики плотностей вероятности}

\begin{figure}[h]
    \centering
    \begin{tikzpicture}[scale=0.9]
        % График нормального распределения (мгновенные значения)
        \begin{scope}[shift={(0,0)}]
            \draw[->, thick] (-4,0) -- (4.5,0) node[right] {$u$, мкВ};
            \draw[->, thick] (0,-0.5) -- (0,5.5) node[above] {$W(u) \times 10^5$, В$^{-1}$};
            
            % Функция Гаусса. Масштаб по x: 1 единица = 10 мкВ. Пик = 1/sqrt(2*pi) * 10 ~ 3.989
            \draw[domain=-3.5:3.5, samples=100, smooth, very thick, blue] plot ({\x}, {3.989*exp(-(\x*10)^2/200)});
            
            \node[below left] at (0,0) {0};
            \draw (1,0.1) -- (1,-0.1) node[below] {10};
            \draw (2,0.1) -- (2,-0.1) node[below] {20};
            \draw (3,0.1) -- (3,-0.1) node[below] {30};
            \draw (-1,0.1) -- (-1,-0.1) node[below] {-10};
            \draw (-2,0.1) -- (-2,-0.1) node[below] {-20};
            \draw (-3,0.1) -- (-3,-0.1) node[below] {-30};
            
            \node[blue, above] at (2, 2) {Нормальное};
        \end{scope}

        % График распределения Релея (огибающая)
        \begin{scope}[shift={(9,0)}]
            \draw[->, thick] (-0.5,0) -- (5,0) node[right] {$E$, мкВ};
            \draw[->, thick] (0,-0.5) -- (0,7) node[above] {$W(E) \times 10^5$, В$^{-1}$};
            
            % Функция Релея. Максимум при x=10 (на графике 1). Пик = 0.1 * exp(-0.5) * 10^5 ~ 6.065
            \draw[domain=0:4.5, samples=100, smooth, very thick, red] plot ({\x}, {10*\x*exp(-(\x*10)^2/200)});
            
            \node[below left] at (0,0) {0};
            \draw (1,0.1) -- (1,-0.1) node[below] {10};
            \draw (2,0.1) -- (2,-0.1) node[below] {20};
            \draw (3,0.1) -- (3,-0.1) node[below] {30};
            \draw (4,0.1) -- (4,-0.1) node[below] {40};
            
            % Линия максимума
            \draw[dashed] (1,0) -- (1, {10*1*exp(-0.5)});
            
            \node[red, right] at (2, 3) {Релея};
        \end{scope}
    \end{tikzpicture}
    \caption{Плотности вероятности мгновенных значений (слева) и огибающей (справа) шума}
\end{figure}